\section{Short labels for prescribed real roots}\label{sec:labels}

A polynomial of exponentially large degree may have exponentially many
real roots, and the derivative signs that specify one root may form an
exponentially long string.  We need to name any prescribed root by only
polynomially many bits.  The following construction replaces that string
by one weighted sum and expresses agreement with the sum by positivity of
an integer polynomial.  It also allows a separate positivity condition at
the chosen root.

\begin{proposition}[Short root labels]\label{prop:root-labels}
Let $P(t)$ and $b_0(t)$ be controlled families of integer univariate
polynomials, indexed by the same input word.  There is a uniformly
controlled family $L_a(t)\in\Z[t]$, with labels $a$ of polynomial binary
length in a polynomial-time specified range, such that whenever
$P\ne0$:
\begin{enumerate}
\item For every label $a$ and every real root $\alpha$ of $P$,
      either $L_a(\alpha)\ge\tfrac12$ or
      $L_a(\alpha)\le-\tfrac12$.
\item For every label $a$ and every real root $\alpha$ of $P$,
      if $L_a(\alpha)>0$, then $b_0(\alpha)>0$.
\item For every real root $\alpha$ of $P$ with $b_0(\alpha)>0$, some
      label $a$ satisfies $L_a(\alpha)>0$ and $L_a(\beta)<0$ at every
      distinct real root $\beta$ of $P$.
\end{enumerate}
The construction is defined and controlled also when $P=0$, but no
root-selection property is asserted in that case.  It uses a fixed
number of applications of the controlled-family rules, independent of
the degrees.  Taking $b_0=1$ permits every real root to be selected.
\end{proposition}

The proof has three ingredients.  First, nonzero values of the required
sign tests at roots lie in a known range bounded away from zero.  Second,
derivative signs distinguish roots, and a modular weighted sum can
separate all their sign strings.  Finally, a rational approximation to
the sign function converts equality of weighted sums into a polynomial
inequality.  We give the quantitative forms of these ingredients before
assembling the labels.

\subsection{Values and derivative signs at algebraic roots}

The value bound must allow nonmonic polynomials and multiple roots.
The product argument below treats both without taking a squarefree part.
Throughout this section, the norm of a polynomial is the sum of the
absolute values of its coefficients.

\begin{lemma}[Bounds for nonzero test values]\label{lem:root-value-bounds}
Let $P\in\Z[t]$ be nonzero, with $\deg P\le d$ and
$\norm{P}\le2^b$, where $d\ge2$ and $b\ge1$ are integers.  If
$g\in\Z[t]$ satisfies $\deg g\le d$ and
$\norm{g}\le2^{b+d^2}$, put
\[
 B_0=b+d^2+d(b+1),\qquad
 B=d(bd+B_0)+bd+2.
\]
At every complex root $\alpha$ of $P$, either $g(\alpha)=0$ or
\[
                    2^{-B}\le |g(\alpha)|\le2^B.
\]
\end{lemma}

\begin{proof}
A nonzero constant polynomial has no roots, so assume that $P$ has actual
degree $m\ge1$ and leading coefficient $a\ne0$.  The elementary root
bound gives $|\alpha|\le1+2^b\le2^{b+1}$, and consequently
$|g(\alpha)|\le2^{B_0}$.

To obtain the lower bound, form the monic integer polynomial and the
integer test polynomial
\[
 \widetilde P(t)=a^{m-1}P(t/a),\qquad
 \widetilde g(t)=a^d g(t/a).
\]
Their integrality follows by inspecting the coefficients; in
$\widetilde P$, the leading coefficient is $1$, and every other
coefficient contains a nonnegative power of $a$.
Multiplication by $t$ on the free abelian group
$\Z[t]/(\widetilde P)$ has the companion matrix of $\widetilde P$.
Over $\C$ it can be triangularized, with eigenvalues $a\alpha$ counted
with their root multiplicities.  Therefore multiplication by
$\widetilde g$ has eigenvalues $a^d g(\alpha)$ with the same
multiplicities.  Its characteristic polynomial is monic with integer
coefficients.  The product of its nonzero eigenvalues, after the zero
factors are removed, is thus, up to sign, a nonzero integer.

Every one of these eigenvalues has modulus at most $2^{bd+B_0}$.
If a specified $g(\alpha)$ is nonzero, there are at most $m-1$ other
nonzero eigenvalues in the product, so
\[
 |a|^d|g(\alpha)|\ge2^{-(m-1)(bd+B_0)}.
\]
Using $|a|\le2^b$ and $m\le d$ gives the claimed lower bound.
\end{proof}

We shall use this lemma for $b_0$ and for
$b_j=P^{(j)}$, $1\le j\le d$.  If $P$ and $b_0$ both have norm at
most $2^b$, these tests have norm at most $2^{b+d^2}$: a derivative
multiplies each coefficient by a falling factorial bounded by
$d^d\le2^{d^2}$.  Zero derivatives cause no difficulty.

The next fact is the uniqueness underlying Thom encodings; see
\cite[Proposition~2.28]{BasuPollackRoy2006}.  Its short proof also explains why repeated
roots are allowed.

\begin{lemma}[Derivative signs distinguish roots]\label{lem:thom-separation}
For a nonzero real polynomial $P$ of degree at most $d$, distinct real
roots have distinct strings
\[
                 \bigl(\sgn P^{(j)}(\alpha)\bigr)_{1\le j\le d}.
\]
\end{lemma}

\begin{proof}
Suppose that two roots $\alpha<\beta$ have the same string, and let $m$
be the actual degree.  The derivative $P^{(m)}$ has constant nonzero
sign on $[\alpha,\beta]$.  Descend from order $m-1$ to order $1$.
If $P^{(j+1)}$ has constant nonzero sign on the interval, then
$P^{(j)}$ is strictly monotone there.  Its endpoint signs agree by
assumption, so they cannot both be zero; they have the same nonzero
sign, which persists throughout the interval.  The induction shows that
$P'$ has constant nonzero sign.  Hence $P$ is strictly monotone between
two of its roots, a contradiction.
\end{proof}

\subsection{Compressing derivative-sign strings}

Fix a nonzero polynomial $P$ of degree at most $d$, with $d\ge2$, and
write $b_j=P^{(j)}$ for $1\le j\le d$.  The derivative signs distinguish
its roots by Lemma~\ref{lem:thom-separation}, but we cannot record all
$d$ signs.  Instead we encode them by a weighted integer sum.

A label is a triple
$a=(\mu,\gamma,S_*)$ in the range
\begin{equation}\label{eq:root-label-range}
 \begin{gathered}
 W_*=2\max(16,d^3)^4,\\
 3\le\mu\le W_*,\qquad
 0\le\gamma<\mu,\qquad
 |S_*|\le dW_*.
 \end{gathered}
\end{equation}
For every such label define integer weights
\[
                    w_j=\gamma^j\bmod\mu,\qquad 0\le w_j<\mu.
\]
At a real root $\alpha$ of a nonzero $P$, the associated weighted sum is
\[
                 S(\alpha)=\sum_{j=1}^d w_j\sgn b_j(\alpha).
\]
This sum is an integer in $[-dW_*,dW_*]$.  The construction below
accepts precisely those roots for which $S(\alpha)=S_*$ and
$b_0(\alpha)>0$.

We claim that some choices of $\mu,\gamma$ make $S$ injective on
the distinct real roots of $P$.  We use polynomial evaluation to
distinguish the strings, the fingerprinting idea underlying
\cite{Schwartz1980}.  The elementary prime-count estimate
in the proof of Proposition~\ref{prop:crt-sign}, with $Z=\max(16,d^3)$,
supplies a prime $\mu$ with
$d^3<\mu\le W_*$.  For two distinct derivative-sign strings
$\sigma,\tau\in\{-1,0,1\}^d$, the polynomial
\[
                       \sum_{j=1}^d(\sigma_j-\tau_j)X^j
\]
is nonzero over $\mathbb F_\mu$: some coefficient belongs to
$\{-2,-1,1,2\}$ and $\mu>2$.  It has at most $d$ zeros in that field.
There are at most $d$ distinct real roots, so fewer than $d^3$ values
of $\gamma$ are forbidden by all pairs together.  Choose a remaining
$\gamma\in\mathbb F_\mu$.  The sums at distinct roots are then distinct
modulo $\mu$, hence distinct as integers.  For a prescribed root
$\alpha$, taking $S_*=S(\alpha)$ selects it uniquely.  Only three short
integers are recorded; the derivative-sign string is never stored.

The label has $O(\log d)$ bits, even though the derivative-sign
string can have length $d$.  We next express the condition
$S(\alpha)=S_*$, together with $b_0(\alpha)>0$, by one polynomial
positivity test.

\subsection{Approximating the sign function}

The algebraic values in Lemma~\ref{lem:root-value-bounds} can be extremely
small, but their binary exponents have short descriptions.  A product
over these exponents gives a rational function that recognizes their
signs with a uniform error.  This sign approximation and the subsequent
clearing of positive denominators are closely related to the method of
Beigel, Reingold, and Spielman~\cite{BeigelReingoldSpielman1995}; we give
the needed construction explicitly.

\begin{lemma}[Rational sign approximation]\label{lem:rational-sign}
For integers $B\ge1$ and $s\ge1$, define
\begin{align*}
 C_{B,s}(z)
   &=\prod_{k=-B}^{B}
       \bigl(2^{\max(0,k)}+2^{\max(0,-k)}z\bigr)^{2s},\\
 A_{B,s}(z)
   &=\frac{C_{B,s}(z)-C_{B,s}(-z)}
           {C_{B,s}(z)+C_{B,s}(-z)}.
\end{align*}
The denominator is at least $1$ for every real $z$, and $A_{B,s}(0)=0$.
For $2^{-B}\le|z|\le2^B$,
\[
                  |A_{B,s}(z)-\sgn z|\le2\cdot3^{-2s}.
\]
Moreover $C_{B,s}$ is an integer polynomial with
\[
 \deg C_{B,s}=2s(2B+1),\qquad
 \norm{C_{B,s}}\le2^{2s(2B+1)(B+1)}.
\]
\end{lemma}

\begin{proof}
Every factor has an even exponent, so $C_{B,s}(z)\ge0$ on $\R$.
For $z\ge0$, each base in $C_{B,s}(z)$ is at least $1$.
Thus $C_{B,s}(z)+C_{B,s}(-z)\ge C_{B,s}(|z|)\ge1$.
The numerator vanishes at zero, and $A_{B,s}$ is odd.

It suffices to consider $2^{-B}\le z\le2^B$.  Choose an integer
$k\in[-B,B]$ for which $2^k$ and $z$ are within a factor of two.
Then
\[
 \left|\frac{2^k-z}{2^k+z}\right|\le\frac13.
\]
All the other factors in the ratio $C_{B,s}(-z)/C_{B,s}(z)$ have
absolute value at most $1$, so this ratio lies in $[0,3^{-2s}]$.
Writing it as $q$, we have
$|A_{B,s}(z)-1|=2q/(1+q)\le2\cdot3^{-2s}$.
The degree is immediate.  Each linear base has coefficient norm at most
$2^{B+1}$, which gives the norm bound by submultiplicativity.
\end{proof}

\subsection{Constructing and verifying the labels}

\begin{proof}[Proof of Proposition~\ref{prop:root-labels}]
Choose common controlled bounds $d\ge2$ and $b\ge1$ such that
\[
 \deg P,\deg b_0\le d,\qquad
 \norm{P},\norm{b_0}\le2^b.
\]
Use the integers $B_0,B$ from Lemma~\ref{lem:root-value-bounds}, and put
$b_j=P^{(j)}$ for $1\le j\le d$.  For a label $a=(\mu,\gamma,S_*)$ in the range
of~\eqref{eq:root-label-range}, use the weights $w_j$ defined there.

It remains to realize this selection by an integer polynomial.  Set
\[
             M=1+dW_*,\qquad s=16M,
             \qquad C=C_{B,s},\qquad A=A_{B,s}.
\]
For every $0\le j\le d$ and every real root $\alpha$ of a nonzero $P$,
Lemmas~\ref{lem:root-value-bounds} and~\ref{lem:rational-sign} give
\begin{equation}\label{eq:root-sign-error}
        |A(b_j(\alpha))-\sgn b_j(\alpha)|
            \le2\cdot3^{-2s}\le\frac1{16M}.
\end{equation}
For the zero test value this holds because $A(0)=0$.
Consequently
\begin{equation}\label{eq:weighted-sign-error}
 \left|\sum_{j=1}^d w_j A(b_j(\alpha))-S(\alpha)\right|
                \le\frac{dW_*}{16M}<\frac1{16}.
\end{equation}

To clear denominators without division, define integer polynomials
\begin{align*}
 D_j(t)&=C(b_j(t))+C(-b_j(t)),\\
 N_j(t)&=C(b_j(t))-C(-b_j(t)),\\
 E(t)&=\prod_{j=0}^d D_j(t),\qquad
 T_j(t)=N_j(t)\prod_{\substack{0\le i\le d\\i\ne j}}D_i(t).
\end{align*}
Thus $E(t)\ge1$ on $\R$ and $T_j(t)=E(t)A(b_j(t))$ there.
We penalize squared disagreement with $S_*$ and with the required
positive sign, and clear the denominators by defining
\begin{equation}\label{eq:root-label-polynomial}
 L_a(t)=E(t)^2
       -4\left(\sum_{j=1}^d w_jT_j(t)-S_*E(t)\right)^2
       -4\bigl(T_0(t)-E(t)\bigr)^2.
\end{equation}
The formula contains only integer additions and multiplications.
For the sign analysis it is useful to write it as
\[
 L_a(t)=E(t)^2\left(1-4\left[
    \left(\sum_{j=1}^d w_j A(b_j(t))-S_*\right)^2
       +\bigl(A(b_0(t))-1\bigr)^2\right]\right).
\]
If $S(\alpha)=S_*$ and $b_0(\alpha)>0$, both absolute differences in
brackets are at most $1/16$.  The parenthetical factor is at least
$1-8/256>1/2$, so $L_a(\alpha)\ge1/2$.
Otherwise, either the integer $S(\alpha)-S_*$ is nonzero, in which
case the first absolute difference is at least $15/16$, or
$\sgn b_0(\alpha)\in\{0,-1\}$, in which case the second is at least
$15/16$.  The factor is then at most $1-4(15/16)^2<-1/2$,
and $L_a(\alpha)\le-1/2$.  This dichotomy holds for every allowed
label, whether or not $\mu$ is prime or the weighted sums are distinct.
Combined with the separating choices already proved to exist, it gives
all three asserted root properties.

Finally, all degree and log-norm bounds, index endpoints, and label
integers have polynomial binary length; their bounds are obtained in
polynomial bit time from the input bounds.  The
weights are computable by modular powering on short integers.
The derivatives $b_j$, the short-index product defining $C$, its
substitutions in $b_j$, the products defining $E,T_j$, and the sums and
squares in~\eqref{eq:root-label-polynomial} are controlled by
Proposition~\ref{prop:controlled-closure}.  In a product omitting index $j$, the omitted
factor is simply replaced by $1$; no denominator is ever inverted in
finite-field evaluation.  The explicit bounds in
Lemma~\ref{lem:rational-sign}, followed by the product and substitution
bounds, retain polynomial binary length.  There are only a fixed number
of nested constructions.  None requires a test that $P\ne0$, so the
same controlled family is defined on the degenerate inputs $P=0$.
\end{proof}
